\chapter{Theoretical Background and Literature Review}
\label{ch:theoretical_background}

This chapter outlines the theoretical principles and literature review underlying the thesis: SOC operational bottlenecks, Tier 1 statistical and machine learning foundations, Tier 2 stateful agent and local LLM architectures, adversarial threats with cryptographic audit mechanisms, and a survey of existing AI Security systems delineating the research gap addressed by SENTINEL.

\section{SOC Operational Realities and Cognitive Bottlenecks}

Section~\ref{ch:introduction} stated the operational consequence; this section states its architectural cause. Security monitoring in modern SOCs aggregates raw telemetry from Next-Generation Firewalls, Intrusion Detection Systems, and Endpoint Detection and Response agents into centralized SIEM platforms~\cite{siem2021}, at volumes far exceeding human analysis capacity~\cite{bigdata2019}. The aggregation itself is not the problem; the problem is that every device feeding it decides by static signature or deterministic rule, so the only lever available for reducing missed detections is a lower threshold, and a lower threshold produces more alerts rather than better ones~\cite{alertfatigue2022}. SOAR platforms sit above this and automate the response, but a playbook can only fire on a situation someone has already anticipated~\cite{soar2025}.

\section{Online Statistical Theory and Tier 1 Machine Learning}
\label{sec:welford}

A wire-speed tier imposes a constraint before it imposes a choice of method: whatever establishes the traffic baseline must do so in constant time and constant memory, because an unbounded stream cannot be re-read and a per-key history cannot be held. That constraint rules out any two-pass formulation of the variance and admits the incremental recurrences of Welford~\cite{welford1962}, in which the sample mean $\mu_k$ and the running sum of squared deviations $S_k$ are carried forward one observation at a time:
\begin{equation}
\label{eq:welford}
    \mu_k = \mu_{k-1} + \frac{x_k - \mu_{k-1}}{k}, \quad S_k = S_{k-1} + (x_k - \mu_{k-1})(x_k - \mu_k)
\end{equation}
The standard deviation $\sigma_k = \sqrt{S_k / (k-1)}$ therefore remains available at every step, and a $Z$-score can be evaluated against a baseline that is never stored. Two properties matter downstream: the baseline adapts as traffic shifts, and it does so without a retention window an attacker could wait out. The classification and caching stages built on top of it are detailed in Section~\ref{sec:tier1_mlgate}.

\section{Local Large Language Models and Tier 2 Agentic Intelligence}
\label{sec:local_llm}

Running the reasoning tier locally is what makes the data-sovereignty property of this design possible, and it is also what makes memory the binding constraint. The arithmetic is unforgiving: eight billion parameters at half precision occupy roughly 16\,GB for the weights alone, before activations and the key-value cache, which already exceeds the 16\,GB card this thesis targets. Quantization resolves this by representing each weight on a coarser integer grid instead of a floating-point one, trading a controlled amount of numerical fidelity for a proportional reduction in footprint~\cite{quantization2021}. In the GGUF Q4\_K\_M format served through \texttt{llama.\allowbreak cpp}, the Foundation-Sec-8B-Instruct model~\cite{foundation_sec_2024, foundation_sec_instruct} occupies 7--8\,GB, which leaves headroom on the same card for the context window and the retrieval index.

To orchestrate reasoning, Tier 2 employs a LangGraph acyclic state graph~\cite{langgraph2024} integrated with Threat Memory. Grounding that reasoning in retrieved evidence rather than in parametric memory~\cite{rag2020} raises a question the security domain answers differently from general question answering: the queries here mix free prose with tokens that must match literally, such as technique identifiers, CVE numbers and port values. Dense vector search~\cite{faiss2019} handles the former and is unreliable on the latter; sparse term matching~\cite{bm25_2009} is the reverse. Because the two produce scores on incomparable scales, they are combined on \textbf{rank} rather than on score, which is the property Reciprocal Rank Fusion exploits~\cite{rrf2009}:
\begin{equation}
\label{eq:rrf}
    RRF\_Score(d) = \sum_{m \in M} \frac{1}{k + r_m(d)}
\end{equation}
retrieving context directly from 433 STIX MITRE ATT\&CK technique entries and NIST SP 800-61r2 guidance for forensic justification.

\section{Adversarial AI Security and Cryptographic Log Audit}
\label{sec:adversarial_theory}

Integrating LLMs introduces log-substrate prompt injection risks~\cite{promptinjection2023}, where attackers inject malicious instructions into User-Agent headers or Syslog strings to hijack AI reasoning~\cite{owasp_llm2025}. The defining property of this vector is that \textbf{log content is authored by the attacker}: any defense based on recognising malicious content is an unending race against novel phrasings. Pandey and Bhujang~\cite{watchtower2026} measure exactly this exposure against an LLM analyst assistant and report that the danger is not where intuition places it: blunt instruction overrides are suppressed entirely, while persona hijacking still suppresses \textbf{68\%} of malicious logs, and context manipulation against a summarisation task succeeds in \textbf{96\%} of attempts undefended and \textbf{38\%} even with constrained output. A content-recognition filter tuned against the obvious class therefore leaves the effective classes open, which is the empirical case for defending the vector structurally instead. To neutralise the vector structurally, SENTINEL applies Delimited Data Encapsulation with random nonce tags isolating raw inputs, combined with deterministic guardrails constraining outputs to three valid actions (\texttt{BLOCK\_IP}, \texttt{LOG}, \texttt{AWAIT\_HITL}).

To protect forensic records, the system uses HMAC-SHA256 signature chaining~\cite{hmac1997} to seal audit records with a local secret key, producing a tamper-evident audit trail. The scope of this guarantee must be stated precisely: because HMAC uses a \textbf{shared symmetric key}, the mechanism establishes \textbf{integrity and tamper-evidence}, but does \textbf{not} constitute non-repudiation in the cryptographic sense---that property requires asymmetric signatures. This boundary is respected throughout the thesis.

\section{Overview of Related Research and Research Gaps}

Artificial intelligence research applied to Security Operations Centers has evolved across three architectural paradigms~\cite{llmsec_survey2025}.

\textbf{Traditional enterprise defense} comprises SIEM platforms such as IBM QRadar and SOAR systems such as Palo Alto Cortex XSOAR~\cite{siem2021, soar2025}. These correlate telemetry with static rules and execute predefined playbooks, so they fail on variants the rules do not anticipate. The cost of that rigidity is measured on the analysts: interviewing SOC practitioners, AlAhmadi \textit{et al.}~\cite{alahmadi2022} find alarm quality to be the dominant complaint, with participants describing the overwhelming majority of alarms as not worth acting on. The strength of this paradigm is very low latency---precisely the property SENTINEL's Tier 1 inherits rather than discards.

\textbf{Retrieval-augmented classification agents} are represented by CyberRAG~\cite{cyberrag2025}, which orchestrates fine-tuned per-family classifiers behind an LLM agent and reports above 94\% accuracy per attack class on SQL injection, XSS, and SSTI benchmarks. Its focus is the quality of classification and reporting for events already selected for analysis; it does not address the prior question of which events deserve that analysis, so no wire-speed tier bounds the arrival rate at the reasoning stage.

\textbf{Multi-agent and governance-aware agents} comprise AutoBnB-RAG~\cite{autobnb2025}, LanG~\cite{lang2026}, and Policy-Guided Threat Hunting~\cite{splunkllm2026}. They should not be flattened into one description, because they differ in exactly the dimension this thesis measures. AutoBnB-RAG evaluates retrieval-augmented multi-agent incident response inside the \textit{Backdoors \& Breaches} tabletop environment---a study of coordination under simulated incidents rather than of an event-processing pipeline. LanG is the closest prior work: it is likewise LangGraph-based with human-in-the-loop checkpoints, reports \textbf{$\approx$21\,ms} inference and a machine-side mean time to detect of \textbf{1.58\,s}, and---contrary to any claim that such systems are undefended---already ships a two-layer guardrail combining regular expressions with a Llama Prompt Guard~2 semantic classifier at \textbf{98.1\%} F1. Policy-Guided Threat Hunting likewise places non-LLM stages first, using a reconstruction autoencoder and a two-layer deep reinforcement learning triage ahead of LLM contextual analysis.

The gap is therefore narrower, and more honest, than an absence of capability in prior work. Each requirement below is met somewhere: cost tiering by Policy-Guided Threat Hunting, prompt-injection defense by LanG, local operation by several. What no single architecture assembles is the \textbf{combination}---(1) resolving the majority of telemetry at the statistical/ML tier under one millisecond, with the offload rate reported as a function of the attack base rate rather than as a constant; (2) stateful reasoning running locally on a 16\,GB GPU under a retrieval-anchoring rule that withholds any technique the retrieved evidence does not support; and (3) prompt-injection defense whose guarantee is \textbf{structural}---derived from where attacker-authored text is placed relative to the instruction boundary, and therefore independent of payload wording, unlike a probabilistic classifier however well it scores---together with HMAC-SHA256 forensic sealing. Table~\ref{tab:related_work_comparison} contrasts SENTINEL against these paradigms.

\begin{table}[htbp]
\centering
\caption{Comparative analysis between SENTINEL and existing security defense paradigms}
\label{tab:related_work_comparison}
\footnotesize
\setlength{\tabcolsep}{4pt}
\begin{tabularx}{\textwidth}{>{\raggedright\arraybackslash}p{2.6cm} >{\raggedright\arraybackslash}X >{\raggedright\arraybackslash}X >{\raggedright\arraybackslash}X}
\toprule
\textbf{Evaluation Metric} & \textbf{SIEM / SOAR (QRadar, XSOAR)} & \textbf{LLM agent systems (CyberRAG, LanG, Policy-Guided)} & \textbf{SENTINEL Architecture (Thesis)} \\
\midrule
\textbf{Core Technology} & Static correlation rules, rigid Python/JS playbooks & Agentic RAG, multi-agent orchestration, DRL triage & Welford $\mathcal{O}(1)$ + LightGBM T1, LangGraph Agent T2 \\
\textbf{Operational Goal} & Centralized log aggregation and playbook response & Automated incident investigation via LLM reasoning & Resolve traffic at the cheap tier, reserve the expensive tier for cases needing reasoning \\
\textbf{Cost Tiering} & One cheap tier, no reasoning & Present in some (DRL pre-triage), absent in others; offload not reported against base rate & \textbf{Two tiers}: cheap first, expensive second; offload reported on two base rates \\
\textbf{Data Sovereignty} & Local enterprise & Local deployment reported; hardware envelope not always stated & 100\% air-gapped on a 16\,GB GPU \\
\textbf{Prompt Inj. Defense} & Not applicable & Probabilistic classifiers (LanG: regex + Prompt Guard 2, 98.1\% F1) & Structural guarantee, independent of payload wording \\
\textbf{Forensic Audit} & Standard text logs (editable post-compromise) & Not reported & HMAC-SHA256 chain, tamper-evident \\
\bottomrule
\end{tabularx}
\end{table}
